\documentclass[12pt]{article}
\usepackage[UTF8]{inputenc}
\usepackage{thaisupp}
\usepackage{geometry}
\geometry{a4paper, margin=2.5cm}
\begin{document}
\begin{center}
{\LARGE\textbf{ฟิสิกส์ ม.ปลาย — Part 1}}\\[0.3em]
{\Large\textbf{Lesson 8: โมเมนตัมและการชน (Extra) (Extra)}}\\[0.5em]
\rule{0.8\linewidth}{0.5pt}
\end{center}
\section*{คำถาม In-Class}
\begin{enumerate}
\item \textbf{วัตถุมวล 4 kg เคลื่อนที่ด้วยความเร็ว 3 m/s มีโมเมนตัมเท่าใด}
\begin{enumerate}
\item ก) 7 kg·m/s
\item ข) 12 kg·m/s
\item ค) 1.33 kg·m/s
\item ง) 0.75 kg·m/s
\end{enumerate}
\item \textbf{วัตถุมวล 2 kg มีโมเมนตัม 16 kg·m/s วัตถุเคลื่อนที่ด้วยความเร็วเท่าใด}
\begin{enumerate}
\item ก) 8 m/s
\item ข) 14 m/s
\item ค) 18 m/s
\item ง) 32 m/s
\end{enumerate}
\item \textbf{วัตถุ A มวล 3 kg วิ่งด้วยความเร็ว 5 m/s ชนกับ B มวล 2 kg ที่อยู่กับที่ แบบไม่ยืดหยุ่น ความเร็วหลังชนเป็นเท่าใด}
\begin{enumerate}
\item ก) 2.14 m/s
\item ข) 3 m/s
\item ค) 4.5 m/s
\item ง) 5 m/s
\end{enumerate}
\item \textbf{วัตถุ A มวล 4 kg วิ่งด้วยความเร็ว 3 m/s ชนกับ B มวล 2 kg ที่อยู่กับที่แบบยืดหยุ่น วัตถุ A หลังชนมีความเร็วเท่าใด}
\begin{enumerate}
\item ก) -1 m/s
\item ข) 0 m/s
\item ค) 1 m/s
\item ง) 2 m/s
\end{enumerate}
\item \textbf{ข้อใดต่อไปนี้เป็นหน่วยเดียวกับโมเมนตัม}
\begin{enumerate}
\item ก) N·s
\item ข) N/m
\item ค) N·s²
\item ง) N/m²
\end{enumerate}
\item \textbf{รถบรรทุกมวล 5000 kg วิ่งด้วยความเร็ว 10 m/s ชนกับรถเก๋งมวล 1500 kg ที่อยู่กับที่ แล้วเคลื่อนที่ติดกันไป ความเร็วหลังชนเป็นเท่าใด}
\begin{enumerate}
\item ก) 3.85 m/s
\item ข) 6.15 m/s
\item ค) 7.69 m/s
\item ง) 10 m/s
\end{enumerate}
\item \textbf{วัตถุสองก้อนมวล 1 kg และ 4 kg วิ่งเข้าหากันด้วยความเร็ว 6 m/s และ 3 m/s ตามลำดับ แบบยืดหยุ่น หลังชนวัตถุ 1 kg มีความเร็วเท่าใด}
\begin{enumerate}
\item ก) -2 m/s
\item ข) -4 m/s
\item ค) -6 m/s
\item ง) -8 m/s
\end{enumerate}
\item \textbf{วัตถุมวล 0.5 kg กระทบผนังด้วยความเร็ว 20 m/s แล้วสะท้อนกลับด้วยความเร็ว 15 m/s impulse มีค่าเท่าใด}
\begin{enumerate}
\item ก) 2.5 N·s
\item ข) 7.5 N·s
\item ค) 17.5 N·s
\item ง) 35 N·s
\end{enumerate}
\item \textbf{กระสุนปืนมวล 0.02 kg ถูกยิงออกด้วยความเร็ว 600 m/s ถ้าปืนมวล 4 kg ปืนจะถอยหลังด้วยความเร็วเท่าใด}
\begin{enumerate}
\item ก) 1.5 m/s
\item ข) 3 m/s
\item ค) 6 m/s
\item ง) 12 m/s
\end{enumerate}
\item \textbf{วัตถุสองก้อนมวลเท่ากัน วิ่งชนกันแบบยืดหยุ่น ก้อนหนึ่งอยู่นิ่ง อีกก้อนวิ่งด้วยความเร็ว v หลังชนจะเกิดอะไรขึ้น}
\begin{enumerate}
\item ก) ก้อนที่วิ่งหยุด ก้อนที่อยู่นิ่งเคลื่อนที่ด้วย v
\item ข) ทั้งสองก้อนเคลื่อนที่ด้วย v/2
\item ค) ก้อนที่วิ่งสะท้อนกลับด้วย v
\item ง) ทั้งสองก้อนหยุดนิ่ง
\end{enumerate}
\item \textbf{วัตถุเคลื่อนที่เป็นวงกลมด้วยอัตราเร็วคงที่ ข้อใดถูกต้อง}
\begin{enumerate}
\item ก) โมเมนตัมคงที่
\item ข) โมเมนตัมเปลี่ยนแปลง
\item ค) แรงลัพธ์เป็นศูนย์
\item ง) ความเร็วคงที่
\end{enumerate}
\item \textbf{ลูกบอลมวล 2 kg ตกจากที่สูง กระทบพื้นด้วยความเร็ว 10 m/s แล้วกระเด้งขึ้นด้วยความเร็ว 6 m/s ถ้าใช้เวลา 0.05 s แรงเฉลี่ยที่พื้นกระทำมีค่าเท่าใด}
\begin{enumerate}
\item ก) 120 N
\item ข) 320 N
\item ค) 480 N
\item ง) 640 N
\end{enumerate}
\item \textbf{วัตถุ A มวล 5 kg วิ่งด้วยความเร็ว 8 m/s ชนกับ B มวล 10 kg ที่อยู่กับที่แบบยืดหยุ่น พลังงานจลน์รวมหลังชนเป็นเท่าใด}
\begin{enumerate}
\item ก) 80 J
\item ข) 160 J
\item ค) 240 J
\item ง) 320 J
\end{enumerate}
\item \textbf{ลูกบอลมวล 0.3 kg ถูกปล่อยให้ตกจากที่สูง กระทบพื้นแล้วสะท้อนขึ้นด้วยความเร็วเท่าเดิม ข้อใดถูกต้อง}
\begin{enumerate}
\item ก) impulse = 0
\item ข) impulse = mv
\item ค) impulse = 2mv
\item ง) impulse = -2mv
\end{enumerate}
\item \textbf{ในการระเบิดของดอกไม้ไฟ ข้อใดต่อไปนี้ถูกต้อง}
\begin{enumerate}
\item ก) โมเมนตัมของระบบไม่คงที่
\item ข) โมเมนตัมของระบบคงที่
\item ค) พลังงานจลน์คงที่
\item ง) พลังงานศักย์คงที่
\end{enumerate}
\item \textbf{วัตถุ A มวล 3 kg วิ่งด้วยความเร็ว 4 m/s ชนกับ B มวล 6 kg ที่วิ่งด้วยความเร็ว 2 m/s ในทิศเดียวกัน แบบไม่ยืดหยุ่น พลังงานจลน์ที่สูญเสียไปเป็นเท่าใด}
\begin{enumerate}
\item ก) 3 J
\item ข) 6 J
\item ค) 9 J
\item ง) 12 J
\end{enumerate}
\item \textbf{ถ้าวัตถุเบามากๆ วิ่งเข้าชนวัตถุหนักมากๆ ที่อยู่กับที่ แบบยืดหยุ่น วัตถุหนักจะเคลื่อนที่ด้วยความเร็วประมาณเท่าใด}
\begin{enumerate}
\item ก) เท่ากับความเร็วเดิม
\item ข) 2 เท่าความเร็วเดิม
\item ค) ครึ่งหนึ่งความเร็วเดิม
\item ง) แทบไม่ขยับ
\end{enumerate}
\item \textbf{ข้อใดต่อไปนี้เป็นปริมาณเวกเตอร์}
\begin{enumerate}
\item ก) พลังงานจลน์
\item ข) งาน
\item ค) โมเมนตัม
\item ง) พลังงานศักย์
\end{enumerate}
\item \textbf{ลูกบอลมวล 0.1 kg วิ่งกระทบผนังด้วยมุม 30° แล้วสะท้อนออกด้วยมุม 30° ด้วยขนาดความเร็วเท่าเดิม 8 m/s impulse ในแนวตั้งฉากผนังมีค่าเท่าใด}
\begin{enumerate}
\item ก) 0.4 N·s
\item ข) 0.8 N·s
\item ค) 1.2 N·s
\item ง) 1.6 N·s
\end{enumerate}
\item \textbf{รถ A มวล 2000 kg วิ่งด้วยความเร็ว 15 m/s ชนติดกับรถ B มวล 3000 kg ที่วิ่งด้วยความเร็ว 10 m/s ในทิศตรงข้าม ความเร็วหลังชนเป็นเท่าใด}
\begin{enumerate}
\item ก) 0 m/s
\item ข) 3 m/s ในทิศรถ A เดิม
\item ค) 3 m/s ในทิศรถ B เดิม
\item ง) 6 m/s ในทิศรถ A เดิม
\end{enumerate}
\item \textbf{อนุภาค A และ B มวลเท่ากัน วิ่งเข้าหากันด้วยความเร็ว 3 m/s และ 6 m/s แบบยืดหยุ่น หลังชนอนุภาค A มีความเร็วเท่าใด}
\begin{enumerate}
\item ก) -3 m/s
\item ข) 3 m/s
\item ค) -6 m/s
\item ง) 6 m/s
\end{enumerate}
\item \textbf{นักระเบิดมวล 60 kg ยืนบนลื่น (ไม่มีแรงเสียดทาน) ยิงระเบิดมวล 0.1 kg ออกไปด้วยความเร็ว 100 m/s นักระเบิดจะเคลื่อนที่ด้วยความเร็วเท่าใด}
\begin{enumerate}
\item ก) 0.6 m/s
\item ข) 0.17 m/s
\item ค) 0.083 m/s
\item ง) 1.0 m/s
\end{enumerate}
\item \textbf{วัตถุมวล 3 kg ถูกแรง F เปลี่ยนแปลงตามเวลา ถ้าโมเมนตัมเปลี่ยนจาก 6 kg·m/s เป็น 24 kg·m/s ในเวลา 3 s แรงเฉลี่ยมีค่าเท่าใด}
\begin{enumerate}
\item ก) 2 N
\item ข) 4 N
\item ค) 6 N
\item ง) 8 N
\end{enumerate}
\item \textbf{ลูกบอลมวล 1 kg กระทบผนังด้วยความเร็ว 12 m/s ในแนวตั้งฉากแล้วสะท้อนกลับด้วยความเร็ว 8 m/s พลังงานจลน์สูญเสียไปเท่าใด}
\begin{enumerate}
\item ก) 24 J
\item ข) 40 J
\item ค) 64 J
\item ง) 88 J
\end{enumerate}
\item \textbf{ลูกบอล A มวล 2 kg วิ่งด้วยความเร็ว 5 m/s ชนกับ B มวล 3 kg ที่อยู่กับที่แบบยืดหยุ่น หลังชน B วิ่งด้วยความเร็วเท่าใด}
\begin{enumerate}
\item ก) 1 m/s
\item ข) 2 m/s
\item ค) 3 m/s
\item ง) 4 m/s
\end{enumerate}
\item \textbf{จักรยานมวล 12 kg วิ่งด้วยความเร็ว 5 m/s คนขี่มวล 60 kg มีโมเมนตัมรวมเท่าใด}
\begin{enumerate}
\item ก) 60 kg·m/s
\item ข) 180 kg·m/s
\item ค) 240 kg·m/s
\item ง) 360 kg·m/s
\end{enumerate}
\item \textbf{สมมติไม่มีแรงภายนอก วัตถุสองก้อนชนกันแล้วสะท้อนออกห่างกัน ข้อใดต่อไปนี้บอกได้ว่าเป็นการชนแบบยืดหยุ่น}
\begin{enumerate}
\item ก) โมเมนตัมคงที่
\item ข) พลังงานจลน์คงที่
\item ค) ทั้งโมเมนตัมและพลังงานจลน์คงที่
\item ง) ความเร็วคงที่
\end{enumerate}
\item \textbf{ลูกบอลมวล 0.5 kg ถูกตีด้วยแรงเฉลี่ย 200 N ในเวลา 0.02 s ความเร็วของลูกบอลหลังถูกตีเป็นเท่าใด}
\begin{enumerate}
\item ก) 4 m/s
\item ข) 8 m/s
\item ค) 12 m/s
\item ง) 16 m/s
\end{enumerate}
\item \textbf{รถบรรทุกมวล 4000 kg วิ่งด้วยความเร็ว 20 m/s ชนกับรถเก๋งมวล 1000 kg ที่อยู่กับที่ ถ้าพลังงานจลน์หลังชนเป็น 160,000 J จงหาความเร็วหลังชน}
\begin{enumerate}
\item ก) 10 m/s
\item ข) 12 m/s
\item ค) 16 m/s
\item ง) 20 m/s
\end{enumerate}
\item \textbf{วัตถุ A มวล 2 kg วิ่งด้วยความเร็ว 10 m/s ชนกับ B มวล 8 kg ที่อยู่กับที่แบบยืดหยุ่น ความเร็วของ A และ B หลังชนเป็นเท่าใด}
\begin{enumerate}
\item ก) A = -6.67 m/s, B = 3.33 m/s
\item ข) A = -5 m/s, B = 2.5 m/s
\item ค) A = -3.33 m/s, B = 6.67 m/s
\item ง) A = -2 m/s, B = 8 m/s
\end{enumerate}
\end{enumerate}
\end{document}